\chapter{Números naturais}\label{ch:g6:wholes}

Tudo em matemática começa contando. Este capítulo revê como os números
naturais são escritos, comparados e combinados --- e examina com atenção a
divisão, a operação que dá ao mesmo tempo um \omterm{def:g3:division:remainder}{quociente} \emph{e} um \omterm{def:g3:division:remainder}{resto}.

\section{Escrever e comparar números naturais}

\begin{definition}[Valor posicional]\label{def:g6:wholes:place}
O nosso jeito de escrever números usa dez algarismos ($0$ a $9$), e o
\emph{valor de um algarismo depende da posição dele}\index{valor posicional}:
em $5\,208$, o algarismo $5$ conta milhares, o $2$ conta centenas, o $0$
conta dezenas (não há nenhuma) e o $8$ conta unidades:
\[
5\,208 = 5 \times 1000 + 2 \times 100 + 0 \times 10 + 8 .
\]
\end{definition}

\begin{omfigure}
\begin{tikzpicture}[scale=0.9]
  \draw (0,0) grid[xstep=2.3, ystep=0.8] (9.2,1.6);
  \node[font=\small] at (1.15,1.2) {milhares};
  \node[font=\small] at (3.45,1.2) {centenas};
  \node[font=\small] at (5.75,1.2) {dezenas};
  \node[font=\small] at (8.05,1.2) {unidades};
  \node[font=\large, omDef] at (1.15,0.4) {$5$};
  \node[font=\large, omDef] at (3.45,0.4) {$2$};
  \node[font=\large, omDef] at (5.75,0.4) {$0$};
  \node[font=\large, omDef] at (8.05,0.4) {$8$};
\end{tikzpicture}

{\small O quadro de \omterm{def:g6:wholes:place}{valor posicional} de $5\,208$. O $0$ importa: sem ele,
o $5$ e o $2$ escorregariam para as colunas erradas e o número se leria
$528$.}
\end{omfigure}

\begin{method}[Comparar dois números naturais]\label{met:g6:wholes:compare}
\begin{enumerate}
  \item O número com mais algarismos é o maior
        ($1\,203 > 989$).
  \item Se tiverem a mesma quantidade de algarismos, compare algarismo por
        algarismo a partir da \emph{esquerda}; a primeira \omterm{ex:g1:subtraction:difference}{diferença}
        decide: $6\,742 > 6\,715$ porque, na casa das dezenas, $4 > 1$.
\end{enumerate}
Os símbolos são $<$ (``menor que'') e $>$ (``maior que''); a ponta fina do
símbolo aponta para o número menor.
\end{method}

\begin{example}\label{ex:g6:wholes:order}
Ordene do menor para o maior: $908$; $89$; $1\,001$; $980$. Primeiro pela
quantidade de algarismos: $89$ (dois algarismos) vem na frente e $1\,001$
(quatro algarismos) no fim. Entre $908$ e $980$: mesmo algarismo das
centenas, $9$; depois, nas dezenas, $0 < 8$, logo $908 < 980$. Resposta
final:
\[
89 < 908 < 980 < 1\,001 .
\]
\end{example}

\begin{omfigure}
\begin{tikzpicture}[scale=1.05]
  \draw[-Stealth] (-0.4,0) -- (10.8,0);
  \foreach \x in {0,1,...,10} \draw (\x,-0.08) -- (\x,0.08);
  \foreach \x/\l in {0/0, 2/200, 4/400, 6/600, 8/800, 10/1000}
    \node[below, font=\small] at (\x,-0.1) {$\l$};
  \fill[omThm] (0.89,0) circle (2pt) node[above, font=\small] {$89$};
  \fill[omDef] (9.08,0) circle (2pt) node[above, font=\small] {$908$};
  \fill[omMeth] (9.8,0) circle (2pt) node[above right=-2pt, font=\small] {$980$};
\end{tikzpicture}

{\small Os números moram numa reta: menor significa mais à esquerda. Cada
traço aqui vale $100$.}
\end{omfigure}

\section{Adição, subtração, multiplicação}

\begin{example}[Contas armadas]\label{ex:g6:wholes:columns}
Para somar ou subtrair, alinhe as unidades sob as unidades e as dezenas
sob as dezenas, e trabalhe da direita para a esquerda, levando o vai-um
quando for preciso:
\[
\begin{array}{r}
4\,5\,7 \\
+\ 2\,8\,6 \\
\hline
7\,4\,3
\end{array}
\qquad\qquad
\begin{array}{r}
6\,0\,3 \\
-\ 1\,4\,7 \\
\hline
4\,5\,6
\end{array}
\]
Na adição: $7 + 6 = 13$, escreva $3$, vai $1$; depois
$5 + 8 + 1 = 14$, escreva $4$, vai $1$; depois $4 + 2 + 1 = 7$. Confira a
subtração somando de volta: $456 + 147 = 603$.
\end{example}

\begin{remark}[Vocabulário]
O resultado de uma adição é uma \emph{soma}\index{soma}; o de uma
subtração, uma \emph{diferença}\index{diferença}; o de uma multiplicação,
um \emph{produto}\index{produto}; o de uma divisão, um
\emph{quociente}\index{quociente}. Os números que estão sendo combinados
são os \emph{termos} (em $+$ e $-$) ou os \emph{fatores} (em $\times$).
\end{remark}

\begin{example}[Cálculo esperto]\label{ex:g6:wholes:smart}
Reordenar termos ou fatores muitas vezes poupa trabalho:
\[
17 + 58 + 3 = (17 + 3) + 58 = 20 + 58 = 78 ,
\]
\[
25 \times 17 \times 4 = (25 \times 4) \times 17 = 100 \times 17 = 1700 .
\]
\end{example}

\section{Divisão}

\begin{theorem}[Divisão euclidiana]\label{thm:g6:wholes:division}
Dados dois números naturais $a$ (o \emph{dividendo}) e $b \neq 0$ (o
\emph{\omterm{def:g5:division:multiple}{divisor}}), existe exatamente uma maneira de escrever
\[
a = b \times q + r
\qquad\text{com}\qquad
0 \leq r < b .
\]
O número $q$ é o \emph{quociente}\index{quociente} e $r$ é o
\emph{resto}\index{resto} da divisão de $a$ por $b$.
\end{theorem}
\admitted

\begin{example}\label{ex:g6:wholes:division}
Vamos dividir $137$ por $4$, passo a passo.
\begin{enumerate}
  \item Quantas vezes o $4$ cabe em $13$? Três vezes
        ($4 \times 3 = 12$), com \omterm{def:g3:division:remainder}{resto} $13 - 12 = 1$.
  \item Baixe o $7$: quantas vezes o $4$ cabe em $17$? Quatro vezes
        ($4 \times 4 = 16$), com \omterm{def:g3:division:remainder}{resto} $1$.
\end{enumerate}
Logo $q = 34$ e $r = 1$:
\[
137 = 4 \times 34 + 1, \qquad\text{e de fato } 0 \leq 1 < 4 .
\]
Conferência: $4 \times 34 = 136$, e $136 + 1 = 137$.
\end{example}

\begin{omfigure}
\begin{tikzpicture}[scale=0.95]
  \draw[-Stealth] (-0.3,0) -- (11.2,0);
  \foreach \i in {0,1,...,34} \draw ({\i*0.32},-0.06) -- ({\i*0.32},0.06);
  \foreach \i in {0,10,20,30}
    \node[below, font=\footnotesize] at ({\i*0.32},-0.08) {$4\times\i$};
  \draw[omDef, very thick] (0,0.14) -- (10.88,0.14);
  \draw[omThm, very thick] (10.88,0.14) -- (10.96,0.14);
  \fill (10.96,0) circle (1.8pt) node[above, font=\small] {$137$};
  \node[omDef, font=\small] at (5.4,0.55) {$34$ grupos completos de $4$};
\end{tikzpicture}

{\small Dividindo $137$ por $4$: cabem $34$ grupos completos de $4$
(chegando a $136$), e sobra $1$ --- o \omterm{def:g3:division:remainder}{resto} é sempre menor que o
\omterm{def:g5:division:multiple}{divisor}.}
\end{omfigure}

\begin{definition}[Divisível]\label{def:g6:wholes:divisible}
Quando o \omterm{def:g3:division:remainder}{resto} é $0$, dizemos que $a$ é \emph{divisível}\index{divisível}
por $b$: por exemplo, $135$ é divisível por $5$, pois
$135 = 5 \times 27 + 0$.
\end{definition}

\begin{method}[Problemas com divisão]\label{met:g6:wholes:problems}
Depois de calcular $a = bq + r$, volte sempre à pergunta:
\begin{enumerate}
  \item ``Quantas caixas / equipes / páginas completas?'' --- a resposta é
        o \omterm{def:g3:division:remainder}{quociente} $q$;
  \item ``Quantos sobram?'' --- a resposta é o \omterm{def:g3:division:remainder}{resto} $r$;
  \item ``Quantas caixas são necessárias para todos?'' --- se $r > 0$, uma
        a \emph{mais} que o \omterm{def:g3:division:remainder}{quociente}: $q + 1$.
\end{enumerate}
\end{method}

\begin{example}\label{ex:g6:wholes:bus}
$137$ alunos vão a um passeio em ônibus de $40$ lugares. Divisão:
$137 = 40 \times 3 + 17$. Três ônibus levam $120$ alunos e sobram $17$
alunos --- que também precisam de ônibus! Logo são necessários $4$
ônibus, mesmo com o \omterm{def:g3:division:remainder}{quociente} sendo $3$.
\end{example}

\section{Exercícios}

\begin{exercise}[$\star$]\label{exo:g6:wholes:1}
Escreva em algarismos: ``três mil e quarenta e sete''; ``vinte mil e
quinhentos''; ``cento e quatro''. Depois escreva $60\,013$ por extenso.
\end{exercise}

\begin{exercise}[$\star$]\label{exo:g6:wholes:2}
No número $47\,209$: qual é o algarismo das centenas? E o das unidades? O
que conta o algarismo $4$? Escreva o número como \omterm{def:g1:addition:def}{soma} de \omterm{def:g5:division:multiple}{múltiplos} de
$1$, $10$, $100$, \dots\ como na
\cref{def:g6:wholes:place}.
\end{exercise}

\begin{exercise}[$\star$]\label{exo:g6:wholes:3}
Copie e complete com $<$ ou $>$:
\[
507 \;?\; 570, \qquad
1\,099 \;?\; 989, \qquad
23\,456 \;?\; 23\,465, \qquad
9\,999 \;?\; 10\,000 .
\]
\end{exercise}

\begin{exercise}[$\star$]\label{exo:g6:wholes:4}
Ordene do menor para o maior: $2\,043$; $978$; $2\,304$; $2\,034$;
$1\,001$.
\end{exercise}

\begin{exercise}[$\star$]\label{exo:g6:wholes:5}
Arme as contas e calcule: $348 + 576$; \quad $805 - 268$; \quad
$307 \times 26$. Confira a subtração com uma adição.
\end{exercise}

\begin{exercise}[$\star$]\label{exo:g6:wholes:6}
Calcule de modo esperto, agrupando termos ou fatores:
\[
38 + 45 + 12 + 5, \qquad
2 \times 83 \times 50, \qquad
4 \times 78 \times 25 .
\]
\end{exercise}

\begin{exercise}[$\star$]\label{exo:g6:wholes:7}
Para cada divisão, dê o \omterm{def:g3:division:remainder}{quociente} e o \omterm{def:g3:division:remainder}{resto}, e escreva a igualdade
$a = b \times q + r$:
\[
95 \div 7, \qquad
230 \div 6, \qquad
504 \div 8 .
\]
\end{exercise}

\begin{exercise}[$\star$]\label{exo:g6:wholes:8}
Sem dividir, explique por que a igualdade $173 = 8 \times 20 + 13$
\emph{não} é a divisão euclidiana de $173$ por $8$, e encontre o
\omterm{def:g3:division:remainder}{quociente} e o \omterm{def:g3:division:remainder}{resto} corretos.
\end{exercise}

\begin{exercise}[$\star\star$]\label{exo:g6:wholes:9}
Os ovos são embalados em caixas de $12$.
\begin{enumerate}
  \item Uma granja tem $2\,000$ ovos. Quantas caixas completas ela
        consegue encher, e quantos ovos sobram?
  \item Uma padaria precisa de $150$ ovos. Quantas caixas ela tem de
        comprar?
\end{enumerate}
\end{exercise}

\begin{exercise}[$\star\star$]\label{exo:g6:wholes:10}
Hoje é terça-feira. Que dia da semana será daqui a $100$ dias? (A semana
tem $7$ dias: divida e pense no \omterm{def:g3:division:remainder}{resto}.)
\end{exercise}

\begin{exercise}[$\star\star$]\label{exo:g6:wholes:11}
Numa divisão por $9$, o \omterm{def:g3:division:remainder}{quociente} é $56$ e o \omterm{def:g3:division:remainder}{resto} é o maior possível.
Qual é o dividendo?
\end{exercise}

\begin{exercise}[$\star\star\star$]\label{exo:g6:wholes:12}
Sou um número de três algarismos. O meu algarismo das centenas é o dobro
do das unidades, o das dezenas é $0$, e a \omterm{def:g1:addition:def}{soma} dos meus algarismos é $9$.
Quem sou eu? (Raciocine passo a passo e depois confira.)
\end{exercise}

\section{Problema: o pequeno Gauss e a soma dos cem primeiros
números}

\begin{problem}[{Problema de fim de semana --- somar $1 + 2 + 3 + \dots +
100$ em dez segundos, e os números triangulares}]
\label{pb:g6:wholes:1}
Uma história famosa: para manter a turma ocupada, um professor pediu aos
alunos que somassem todos os números de $1$ a $100$. Um menino escreveu
$5\,050$ na lousinha quase em seguida --- Carl Friedrich Gauss, com nove
anos, que se tornaria um dos maiores matemáticos de todos os tempos. O
segredo dele foi apenas \emph{cálculo esperto}
(\cref{ex:g6:wholes:smart}): agrupar os termos de modo que cada grupo
fique fácil. Este problema redescobre o truque dele, faz com que funcione
em qualquer situação e mostra como essas \omterm{def:g1:addition:def}{somas} aparecem no dia a dia.

\medskip
\noindent\textbf{Parte I --- Os pares de Gauss.}
\begin{enumerate}
  \item Calcule $1 + 2 + 3 + 4 + 5$ agrupando os termos com esperteza.
  \item Para $1 + 2 + 3 + \dots + 10$: junte o primeiro termo com o
        último ($1 + 10$), o segundo com o penúltimo ($2 + 9$), e assim
        por diante. Quanto vale cada par? Quantos pares há? Termine a
        conta.
  \item O mesmo método para $1 + 2 + \dots + 20$: quanto vale cada par,
        quantos pares, que total?
  \item Agora a \omterm{def:g1:addition:def}{soma} do professor, $1 + 2 + \dots + 100$: explique por
        que os pares $1 + 100$, $2 + 99$, $3 + 98, \dots$ valem todos o
        mesmo, e por que há exatamente $50$ deles, sendo o último
        $50 + 51$.
  \item Escreva a resposta de Gauss como uma única multiplicação e
        calcule-a.
\end{enumerate}

\medskip
\noindent\textbf{Parte II --- Um truque que nunca falha.}
\begin{enumerate}[resume]
  \item Experimente o emparelhamento em $1 + 2 + \dots + 9$: desta vez um
        número fica sem par. Qual? Emparelhe em torno dele e calcule a
        \omterm{def:g1:addition:def}{soma}.
  \item Eis uma versão do truque que funciona quer a quantidade de termos
        seja par, quer seja \omterm{def:g3:numbers:evenodd}{ímpar}: escreva a \omterm{def:g1:addition:def}{soma} duas vezes, uma para a
        frente e outra de trás para a frente, uma linha embaixo da
        outra,
        \[
        \begin{array}{ccccccc}
        1 & + & 2 & + & \dots & + & 9 \\
        9 & + & 8 & + & \dots & + & 1
        \end{array}
        \]
        e some \emph{coluna por coluna}. Quanto vale cada coluna? Quantas
        colunas há? Explique por que a \omterm{def:g1:addition:def}{soma} das duas linhas é
        $9 \times 10 = 90$ e conclua que $1 + 2 + \dots + 9 = 45$ --- a
        mesma resposta da questão 6.
  \item Use o truque das duas linhas para calcular $1 + 2 + \dots + 50$
        e depois $1 + 2 + \dots + 200$.
  \item Calcule a \omterm{def:g1:addition:def}{soma} dos cem primeiros números \emph{pares},
        $2 + 4 + 6 + \dots + 200$. (Compare cada termo com um termo da
        \omterm{def:g1:addition:def}{soma} de Gauss: não é preciso nenhum emparelhamento novo.)
  \item Calcule $101 + 102 + \dots + 200$, usando duas \omterm{def:g1:addition:def}{somas} que você já
        conhece e uma subtração.
\end{enumerate}

\medskip
\noindent\textbf{Parte III --- \omterm{pb:g6:wholes:1}{Números triangulares} por toda parte.}
Os totais $1$, $3$, $6$, $10$, $15, \dots$ das \omterm{def:g1:addition:def}{somas}
$1 + 2 + \dots + n$ chamam-se \emph{números
triangulares}\index{números triangulares}, porque contam objetos
empilhados em triângulo.
\begin{enumerate}[resume]
  \item Oito amigos se encontram, e cada um aperta a mão de cada um dos
        outros exatamente uma vez. Explique por que o número de apertos
        de mão é $7 + 6 + 5 + 4 + 3 + 2 + 1$ e calcule-o.
  \item Uma quitandeira monta uma \omterm{def:g5:solids:def}{pirâmide} triangular de latas: $12$
        latas na fileira de baixo, $11$ na seguinte, e assim por diante
        até uma única lata no topo. De quantas latas ela precisa? (Use o
        truque da questão 7.)
  \item Desenhe a \omterm{def:g1:addition:def}{soma} $1 + 2 + 3 + 4$ como uma escada de quadradinhos:
        um quadradinho na primeira coluna, dois na segunda, três na
        terceira, quatro na quarta. Mostre no seu desenho que
        \emph{duas} escadas dessas, uma delas virada de cabeça para
        baixo, se encaixam exatamente num retângulo de $4$ fileiras de
        $5$ quadradinhos --- e explique por que isso é o truque das duas
        linhas da questão 7, desenhado em vez de escrito.
  \item Um relógio bate $1$ vez à uma hora, $2$ vezes às duas horas,
        \dots, $12$ vezes às doze horas. Quantas badaladas são essas ao
        longo das doze horas? E ao longo de um dia inteiro?
  \item Uma última para Gauss: $5\,050$ segundos --- é mais ou menos que
        uma hora e meia? Converta em horas, minutos e segundos, usando
        duas vezes a divisão euclidiana
        (\cref{thm:g6:wholes:division}).
\end{enumerate}
\end{problem}
